\section{Pengenalan Pengujian (Testing) Aplikasi Web dan Tujuannya}

\textbf{Pengujian} (\textit{testing}) adalah proses sistematis untuk memverifikasi bahwa aplikasi web berfungsi sesuai harapan, bebas bug, dan memenuhi standar kualitas. Pengujian bukan langkah terakhir --- idealnya diintegrasikan di setiap tahap pengembangan (\textit{shift-left testing}). Tanpa pengujian, bug ditemukan oleh pengguna di produksi, yang jauh lebih mahal untuk diperbaiki dibandingkan saat pengembangan \cite{mdnToolsTesting}.

Tujuan pengujian meliputi: (1) \textbf{verifikasi fungsionalitas} --- memastikan fitur bekerja sesuai spesifikasi; (2) \textbf{deteksi bug} --- menemukan kesalahan sebelum deployment; (3) \textbf{keamanan} --- memastikan tidak ada celah keamanan; (4) \textbf{performa} --- memastikan aplikasi cepat dan responsif; (5) \textbf{kompatibilitas} --- memastikan berjalan di berbagai browser dan perangkat. Pengujian yang baik memberikan kepercayaan untuk melakukan perubahan dan refactoring tanpa merusak fitur yang ada \cite{mdnToolsTesting}.

\begin{table}[ht]
\centering
\begin{tabular}{|P{3cm}|P{4.5cm}|P{5cm}|}
\hline
\textbf{Tujuan} & \textbf{Deskripsi} & \textbf{Contoh} \\
\hline
Verifikasi & Fitur sesuai spesifikasi & Form registrasi menyimpan data \\
\hline
Deteksi bug & Temukan error sebelum rilis & Tombol submit tidak merespons \\
\hline
Keamanan & Tidak ada celah serangan & SQL injection via form login \\
\hline
Performa & Aplikasi cepat dan responsif & Halaman load $<$ 3 detik \\
\hline
Kompatibilitas & Berjalan di semua browser & Layout benar di Chrome/Firefox \\
\hline
\end{tabular}
\caption{Tujuan utama pengujian aplikasi web}
\end{table}

